\begin{samepage}
\section{Nodal length at the wavelength scale}\label{sec:length}

We first state the planar estimate that relates nodal length to
growth. The disks in this section are Euclidean.

\begin{theorem}[Roy-Fortin]\label{thm:local-length}
There are constants $\eps_0>0$ and $C>0$ with the following property.
Let $F$ be a nonzero real solution of
\[
 \Delta F+qF=0\quad\text{on }\D_3,
 \qquad q\in C^\infty(\D_3),\quad \norm{q}_\infty<\eps_0.
\]
If
\[
 \beta=\log\frac{\norm{F}_{L^\infty(\D_{5/2})}}
                         {\norm{F}_{L^\infty(\D_{1/4})}},
\]
then
\[
 \mathcal H^1(Z_F\cap\D_{1/60})\le C\max\{1,\beta\}.
\]
\end{theorem}
\end{samepage}

This is \cite[Theorem~2.1.1]{RoyFortin2015}. The potential need not
have a constant sign. We use the following consequence, keeping
the norm and length scales explicit.

\begin{lemma}[Nodal length from local $L^2$ growth]\label{lem:wavelength-length}
Fix the coefficient bounds and coordinate neighborhood from
Section~\ref{sec:subdivision}. There is an $\eps>0$ such that,
for every admissible square $I$ of side $s$ with $ks\le\eps$,
\begin{equation}\label{eq:wavelength-length}
 \mathcal H^1(I\cap Z_u)\le Cs\bigl(1+E(I,0)\bigr).
\end{equation}
The constants are independent of $u$, $k$, and $s$.
\end{lemma}

\begin{proof}
Let $x_I$ be the center of $I$, and put
\[
 F(z)=u(x_I+60sz),\qquad q(z)=(60ks)^2p(x_I+60sz).
\]
The disk of physical radius $180s$ is available because $R=1000$.
Choose $\eps$ so small that
$(60\eps)^2\norm{p}_\infty<\eps_0$ uniformly over the coordinate
regions in use. Then $F$ satisfies Theorem~\ref{thm:local-length}.

An interior $H^2$ estimate for $\Delta F=-qF$, followed by the
two-dimensional Sobolev embedding, gives
\[
 \norm{F}_{L^\infty(\D_{5/2})}
 \le C\norm{F}_{L^2(\D_{11/4})}.
\]
Only the uniform bound on $q$ is needed here. On the smaller disk,
\[
 \norm{F}_{L^2(\D_{1/4})}\le
 |\D_{1/4}|^{1/2}\norm{F}_{L^\infty(\D_{1/4})}.
\]
Consequently
\[
 \beta\le C+
 \log\frac{\norm{F}_{L^2(\D_{11/4})}}
               {\norm{F}_{L^2(\D_{1/4})}}.
\]
The same $L^2$ version of the local length bound is recorded directly
in \cite[Theorem~3.1.1]{RoyFortinLq}.

In physical coordinates the two norm radii are $165s$ and $15s$.
Their common change-of-variables factor cancels, and nesting gives
\[
 \log\frac{\norm{F}_{L^2(\D_{11/4})}}
               {\norm{F}_{L^2(\D_{1/4})}}
 =\log\frac{\norm{u}_{L^2(\D_{165s}(x_I))}}
                {\norm{u}_{L^2(\D_{15s}(x_I))}}
 \le E(I,0).
\]
All these norms are nonzero by Proposition~\ref{prop:initial-growth}.
The length disk becomes $\D_s(x_I)$ and contains the closed square
$I$, whose corner distance is $s/\sqrt2$. Hausdorff length scales
by $60s$. Theorem~\ref{thm:local-length} therefore yields
\eqref{eq:wavelength-length}.
\end{proof}

\begin{proof}[Proof of Theorem~\ref{thm:main}]
Choose finitely many smooth isothermal neighborhoods and root
squares $Q_0$ whose interiors cover $M$. Make their sides
$s_0\le1$ small enough that each closed disk of radius $2Rs_0$
about a root center is contained in its coordinate neighborhood.
The conformal factors and their first derivatives have uniform
bounds on the finitely many regions used below, and the conformal
factors are bounded away from zero there. Proposition~\ref{prop:initial-growth}
gives a uniform initial bound $N(Q_0)\le C_{\mathrm{init}}$.

Choose the subdivision constants from
Proposition~\ref{prop:subdivision}, and then the small
$\eps$ of Lemma~\ref{lem:wavelength-length}. Write
$k_{\min}=\sqrt{\lambda_1(M,g)}>0$. A single permissible lower
scale bound for all roots is
\[
 a_0=\min\left\{\frac{\eps}{A},\,
           k_{\min}\min_{Q_0}s(Q_0)\right\}>0.
\]
Now fix the descent and persistence constants with this $a_0$.
This order of choices makes every constant depend only on the
fixed metric and cover.

For a given eigenfunction, subdivide each root to the first level
at which $ks\le\eps$, where $s$ is that level's common side.
If this level is below the root, its parent has $kAs>\eps$, so
$ks>\eps/A$. If it is the root, then $ks\ge k_{\min}s_0$.
In both cases, every visited square satisfies $K_Q\ge a_0$.
Theorem~\ref{thm:mean-growth} therefore gives
\[
 \av_{I\text{ terminal in }Q_0}E(I,0)\le Cks.
\]
There are $(s_0/s)^2$ terminal squares. Applying
Lemma~\ref{lem:wavelength-length} and summing their lengths yields
\begin{align*}
 \mathcal H^1(Q_0\cap Z_u)
 &\le Cs\left(\frac{s_0}{s}\right)^2(1+Cks)\\
 &\le Cs_0^2(1/s+k)
 \le Cs_0^2(a_0^{-1}+1)k.
\end{align*}
Subadditivity suffices even if a nodal arc lies on a grid edge.

On each convex coordinate square, the inverse coordinate map is
Lipschitz into $(M,d_g)$ with constant at most
$\sqrt{\sup p}$: the coordinate line segment is an admissible
path for the Riemannian distance. Hence the preceding Euclidean
Hausdorff-length bound gives the same bound for $\mathcal H_g^1$
up to a fixed factor. Sum over the finite root cover to obtain
$\mathcal H_g^1(Z_u)\le C(M,g)k$.
The coordinate choices and all constants were independent of
the eigenfunction and of its eigenspace. Since $k=\sqrt\lambda$,
this proves the theorem.
\end{proof}
